\scanpage{102}
\mspage{101}
% Upper-left margin: thin grey curved line; origin unassigned.
\margin{\figure{\begin{tikzpicture}\draw[gray,line width=.2pt] (0,0) .. controls (1,-.13) and (2.2,-.19) .. (3.2,-.14);\end{tikzpicture}}}
% Upper-right margin: two separate diffuse brown-grey traces with darker centers; no textual value or maker assigned.
\margin{\figure{\begin{tikzpicture}
\fill[brown!30!gray,opacity=.45] (-.018,.02) .. controls (.03,.01) and (.04,-.08) .. (.027,-.17) .. controls (-.01,-.20) and (-.04,-.09) .. cycle;
\fill[black!65] (-.012,-.145) .. controls (.015,-.157) and (.023,-.172) .. (.004,-.181) .. controls (-.018,-.178) and (-.024,-.159) .. cycle;
\fill[brown!30!gray,opacity=.45] (.385,-.66) .. controls (.42,-.72) and (.43,-.83) .. (.408,-.96) .. controls (.39,-1.07) and (.36,-1.14) .. (.348,-1.18) .. controls (.326,-1.19) and (.341,-1.08) .. (.365,-.96) .. controls (.39,-.83) and (.36,-.72) .. cycle;
\fill[black!65] (.372,-.90) ellipse (.015 and .018);
\fill[black!65] (.371,-1.01) .. controls (.391,-1.035) and (.376,-1.06) .. (.356,-1.063) .. controls (.348,-1.045) and (.36,-1.02) .. cycle;
\fill[black!45] (.350,-1.10) .. controls (.373,-1.085) and (.371,-1.13) .. (.341,-1.16) .. controls (.335,-1.145) and (.341,-1.12) .. cycle;
\end{tikzpicture}}}

Let us now examine the case where one of the cycles $\underline{a}$ or $\underline{b}$ (for example $\underline{b}$) belongs to the Albanese kernel $Z_{\handwritten{\ell}}(A)$ ; by linearity, we reduce to the case where $\underline{b}$ is of the form $\underline{c}_u-\underline{c}$, with $\underline{c}\handwritten{\in}Z^{s}_{\unsure{o|0}}(A)_{\bar{\underline{k}}}$, and $u\handwritten{\in}{}^{s}A_{\bar{\underline{k}}}$. By the invariance of the symbol $j_{\handwritten{o}}$ under the translation $\handwritten{\tau_u}$, the two sides of (3) are respectively equal to $j_{\handwritten{o}}(Y_{\underline{a}},\underline{c})$ and to $j_{\handwritten{o}}(Y^{-}_{\underline{c}},\underline{a})$\inserted{\handwritten{, setting $Y=X_u-X$.}} Since we have $X\handwritten{\in}D_a(A)$, we have $Y\handwritten{\sim}0$, and we are reduced to the case already treated.
% The handwritten insertion begins with comma + en at the right end of line 6;
% posant Y=X_u-X. continues in the left margin at line 7, before Puisqu'on.

Suppose now that $X$, $\underline{a}$ and $\underline{b}$ are arbitrary. We can, by modifying $\underline{b}$ by the addition of an element of $Z^{\handwritten{s}}_{\handwritten{\ell}}(A)_{\handwritten{\bar{\underline{k}}}}$, reduce to the case where the cycle $\underline{a}^{-}\handwritten{\otimes}\underline{b}$ is disjoint from $X$. The symbols $X_{\underline{a}}(\underline{b})$ and $X^{-}_{\underline{b}}(\underline{a})$ are then defined, and we have
\[
X_{\underline{a}}(\underline{b})=X^{-}_{\underline{b}}(\underline{a}),\quad\text{whence}\quad v(X_{\underline{a}}(\underline{b}))=v(X^{-}_{\underline{b}}(\underline{a})).
\]

On the other hand, let $a$ (resp. $b$) be any one of the components of $\underline{a}$ (resp. $\underline{b}$). Then the symbols $i(X_a,b)$ and $i(X^{-}_b,a)$ are defined ; the map $\handwritten{\unsure{\zeta|\xi}}:A\handwritten{\longrightarrow}A$ defined by $\handwritten{\zeta}(x)=a+b-x$ being \uline{p}-isomorphic at $b^{\handwritten{\unsure{\circ|o|0}}}$, with value $a^{\handwritten{\circ}}$, we have, by the invariance property of the symbol $i$,
\[
i(X_a,b)=i(X^{-}_b,a).
\]

It follows, by linearity,
\[
i(X_{\underline{a}},\underline{b})=i(X^{-}_{\underline{b}},\underline{a}).
\]
% Left margin beside the second i display: gray stem and dark filled lobes; origin unassigned.
\margin{\figure{\begin{tikzpicture}[x=1cm,y=1cm]
\fill[black!35] (0,0) .. controls (-.004,-.04) and (.006,-.08) .. (.018,-.11) .. controls (.023,-.14) and (.020,-.17) .. (.026,-.20) .. controls (.040,-.24) and (.045,-.28) .. (.037,-.30) .. controls (.018,-.31) and (.009,-.27) .. (.013,-.24) .. controls (.020,-.21) and (.022,-.20) .. (.012,-.18) .. controls (.005,-.15) and (.011,-.13) .. (.008,-.10) .. controls (-.003,-.06) and (-.006,-.02) .. (0,0);
\fill[black] (.019,-.12) .. controls (.030,-.14) and (.022,-.17) .. (.026,-.18) .. controls (.008,-.17) and (.008,-.15) .. (.019,-.12);
\fill[black] (.024,-.22) .. controls (.034,-.24) and (.044,-.28) .. (.037,-.30) .. controls (.019,-.31) and (.011,-.28) .. (.013,-.25) .. controls (.014,-.23) and (.019,-.22) .. (.024,-.22);
\end{tikzpicture}}}

Taking into account the preceding formulas (2), we deduce (3).

Note finally \unsure{.}that $j_{\handwritten{o}}(X,\underline{a})$ vanishes when all the components of $A$ belong to the group ${}^{s}A_{\bar{\underline{k}},\handwritten{o}}$. Indeed, let $f$ be a function on $V$, defined over $\underline{k}$, such that

% archive sheet number: 101
